\section{Лекция 16}
\subsection{Свойства криволинейного интеграла первого рода}
\begin{statement} [Необходимое условие существования] $\\$
    Если существует интеграл
    \[\int\limits_{\gamma}f\ ds\]
    то $f$ ограничена на $\gamma$.
\end{statement}
\begin{proof}
    По необходимому условию для интеграла Римана-Стилтьеса.
\end{proof}
\begin{reminder}
    Если $f\in \mathcal{C}([a,b]),\ g$ --- функция ограниченной вариации, то существует интеграл
    \[\int\limits_{a}^{b}f(x)\ d(g(x))\]
\end{reminder}
\begin{statement}[Достаточное условие существования] $\\$
    Если $f\in \mathcal{C}(\phi([\alpha,\beta]))$ и $\gamma$ спрямляема, то существует интеграл
    \[\int\limits_{\gamma}f\ ds\]
\end{statement}
\begin{proof}
    Поскольку кривая $\gamma$ спрямляема, то функция длины дуги $s(t)$ определена, непрерывна и монотонно возрастает на $[a,b]$. Известно, что $f\in \mathcal{C}(\phi([\alpha,\beta]))$, значит $f(\phi(t))$ непрерывна на $[\alpha,\beta]$. Монотонная функция $s(t)$ является функцией ограниченной вариации. Тогда по свойству для интеграла Римана-Стилтьеса интеграл существует.
\end{proof}
\begin{statement} [Линейность] $\\$
    Пусть $c_1,c_2\in \R$ и существуют интегралы
    \[\int\limits_{\gamma}f\ ds,\ \ \int\limits_{\gamma}g\ ds\]
    Тогда существует интеграл
    \[\int\limits_{\gamma}(c_1\cdot f+c_2\cdot g)\ ds=c_1\cdot \int\limits_{\gamma}f\ ds+c_2\cdot \int\limits_{\gamma}g\ ds\]
\end{statement}
\begin{proof}
    По линейности интеграла Римана-Стилтьеса.
\end{proof}
\begin{reminder}
    Пусть $f$ --- функция ограниченной вариации, $g\in \mathcal{D}([a,b]),\\
    g'\in \mathcal{R}([a,b])$. Тогда интеграл Римана-Стилтьеса сводится к интегралу Римана:
    \[(S)\int\limits_{a}^{b}f(x)\ d(g(x))=(R)\int\limits_{a}^{b}f(x)\cdot g'(x)\ dx\]
\end{reminder}
\begin{statement}
    Пусть $\phi: [\alpha,\beta]\to \R^k$ --- параметризация кривой $\gamma$ и $\phi\in \mathcal{C}^1$. Тогда
    \[\int\limits_{\gamma}f\ ds=\int\limits_{\alpha}^{\beta}f(\phi)\cdot \sqrt{(x_1'(t))^2+\dots+(x_k'(t))^2}\ dt\]
\end{statement}
\begin{proof}
    Сведем интеграл Римана-Стилтьеса к интегралу Римана:
    \begin{multline*}
        \int\limits_{\gamma}f\ ds=\int\limits_{\alpha}^{\beta}f(\phi(t))\ ds(t)=\int\limits_{\alpha}^{\beta}f(\phi(t))\cdot \|\phi'(t)\|\ dt=\\
        =\int\limits_{\alpha}^{\beta}f(\phi(t))\cdot \sqrt{(x_1'(t))^2+\dots+(x_k'(t))^2}\ dt
    \end{multline*}
\end{proof}
\begin{statement}
    Интеграл
    \[\int\limits_{\gamma}f\ ds\]
    не зависит от ориентации.
\end{statement}
\begin{proof}
    Рассмотрим интегральную сумму
    \[\sigma_{s}(f\circ\phi, T, H)=\sum\limits_{i=1}^{n}f(\xi_i)\Delta s_i\]
    При смене ориентации слагаемые будут записанны в другом порядке, это ни как не повлияет на интеграл.
\end{proof}
\begin{statement}
    Пусть $f\geq g$ на $\gamma$ и существуют интегралы
    \[\int\limits_{\gamma}f\ ds,\ \ \int\limits_{\gamma}g\ ds\]
    Тогда
    \[\int\limits_{\gamma}f\ ds\geq\int\limits_{\gamma}g\ ds\]
\end{statement}
\begin{proof} Заметим, что
    \[\sigma_s(f\circ\phi,T,H)\geq\sigma_s(g\circ\phi,T,H)\]
    Устремим $d(T)\to 0$ и получим утверждение.
\end{proof}
\begin{statement}[Аддитивность] $\\$
    Если $\gamma=\gamma_1\sqcup\gamma_2$ и существуют интегралы
    \[\int\limits_{\gamma}f\ ds,\ \ \int\limits_{\gamma_1}f\ ds,\ \ \int\limits_{\gamma_2}f\ ds\]
    Тогда
    \[\int\limits_{\gamma}f\ ds=\int\limits_{\gamma_1}f\ ds+\int\limits_{\gamma_2}f\ ds\]
\end{statement}
\begin{proof}
    Пусть $T_1$ --- разбиение $\gamma_1$, а $T_2$ --- разбиение $\gamma_2$. Рассмотрим разбиение $T$, которое содержит точку стыка $\gamma_1$ и $\gamma_2$ и совпадает с $T_1$ и $T_2$ на $\gamma_1$ и $\gamma_2$ соответсвенно. Тогда
    \[\sigma_s(f\circ \phi, T,H)=\sigma_s(f\circ \phi, T_1,H)+\sigma_s(f\circ \phi, T_2,H)\]
    Устремим $d(T)\to 0$ и получим утверждение теоремы.
\end{proof}
\begin{statement}
    Пусть $f\geq 0$ на $\gamma=\gamma_1\sqcup\gamma_2$. Тогда
    \[\int\limits_{\gamma}f\ ds\geq \int\limits_{\gamma_1}f\ ds\]
\end{statement}
\begin{proof}
    Следует из аддитивности и неотрицательности.
\end{proof}
\begin{statement}
    Пусть $\phi$ --- параметризация кривой $\gamma$,\ $f$ непрерывна на $\phi([\alpha,\beta])$ и $\gamma$ --- спрямляема. Тогда
    \[\int\limits_{\gamma}f\ ds\leq M\cdot |\gamma|\]
    где $M=\sup\limits_{\gamma}f$.
\end{statement}
\begin{proof}
    \[\sigma_s(f\circ\phi,T,H)=\sum\limits_{i=1}^{n}f(\xi_i)\Delta s_i\leq M\cdot \sum\limits_{i=1}^{n}\Delta s_i=M\cdot |\gamma|\]
    Устремим $d(T)\to 0$ и получим утверждение теоремы.
\end{proof}
\subsection{Криволинейный интеграл второго рода}
\begin{definition}
    Пусть $\gamma$ --- ориентированная кривая с параметризацией 
    \[\phi:[\alpha,\beta]\to \R^k,\ \ \phi(t)=(x_1(t),\dots,x_k(t))\]
    Пусть $f_1,\dots,f_k$ --- функции, определенные на $\gamma$ такие, что для каждого\\
    $i=1,\dots,k$ существует интеграл
    \[\int\limits_{\alpha}^{\beta}f_i(x_1(t),\dots,x_k(t))\ d(x_i(t))\]
    Тогда
    \[\int\limits_{\gamma}f_1\ dx_1+f_2\ dx_2+\dots+f_k\ dx_k := \sum\limits_{i=1}^{k}\int\limits_{\alpha}^{\beta}f_i(\phi(t))\ d(x_i(t))\]
    называется криволинейным интегралом второго рода.
\end{definition}
\begin{definition}[Альтернативное определение криволинейного интеграла второго рода]\tab
    \begin{enumerate}
        \item Пусть $\gamma$ --- кривая с параметризацией 
        \[\phi:[\alpha,\beta]\to \R^k,\ \ \phi(t)=(x_1(t),\dots,x_k(t))\]
        Пусть $f_1,\dots,f_k$ --- функции, определенные на $\gamma$. Рассмотрим $T=\{t_i\}$ и\\
        $H=\{\xi_i\}$ --- произвольные разбиение и разметка отрезка $[\alpha,\beta]$.\
        \item Обозначим $\Delta_{ij}=x_i(t_j)-x_i(t_{j-1})$. Интегральной суммой назовем
        \[\sigma(f,T,H)=\sum\limits_{j=1}^{n}\Big(f_1(\phi(\xi_j))\Delta x_{1j}+f_2(\phi(\xi_j))\Delta x_{2j}+\dots+f_k(\phi(\xi_j))\Delta x_{kj}\Big)\]
        \item Наконец, если
        \[\forall \epsilon>0\ \exists\ \delta>0: \forall T,\ d(T)<\delta: |\sigma(f,T,H)-I|<\epsilon\]
        то 
        \[I=\int\limits_{\gamma}f_1\ dx_1+f_2\ dx_2+\dots+f_k\ dx_k\]
        называется криволинейным интегралом второго рода.
    \end{enumerate}
\end{definition}
\subsection{Свойства криволинейного интеграла второго рода}
\begin{statement} [Достаточное условие существования] $\\$
    Пусть $f_1,\dots,f_k$ --- непрерывны на спрямляемой кривой $\gamma$. Тогда существует интеграл
    \[\int\limits_{\gamma}f_1\ dx_1+f_2\ dx_2+\dots+f_k\ dx_k\]
\end{statement}
\begin{proof}
    Поскольку $\gamma$ спрямляема, то каждая $x_i(t)$ имеет ограниченную вариацию. При этом $f_i(\phi(t))$ непрерына на $[\alpha,\beta]$. Тогда по свойству для интеграла Римана-Стилтьеса интеграл существует.
\end{proof}
\begin{statement}
    Пусть $c_1,c_2\in \R$
    %\[\vec{F}=(f_1,\dots,f_k),\ \ \vec{G}=(g_1,\dots,g_k)\] 
    и существуют интегралы
    \[\int\limits_{\gamma}f_1\ dx_1+f_2\ dx_2+\dots+f_k\ dx_k,\ \ \int\limits_{\gamma}g_1\ dx_1+g_2\ dx_2+\dots+g_k\ dx_k\]
    Тогда существует интеграл
    \begin{multline*}
        \int\limits_{\gamma}(c_1\cdot f_1+c_2\cdot g_1)\ dx_1+\dots+(c_1\cdot f_k+c_2\cdot g_k)\ dx_k=\\
        =c_1\cdot \int\limits_{\gamma}f_1\ dx_1+\dots+f_k\ dx_k+c_2\cdot \int\limits_{\gamma}g_1\ dx_1+\dots+g_k\ dx_k
    \end{multline*}
\end{statement}
\begin{proof}
    По линейности интеграла Римана-Стилтьеса для каждой координаты.
\end{proof}
\begin{statement}
    Пусть $\gamma$ --- спрямляемая кривая, $\phi\in \mathcal{C}^1([\alpha,\beta])$. Тогда
    \[\int\limits_{\gamma}f_1\ dx_1+\dots+f_k\ dx_k=\int\limits_{\alpha}^{\beta}\Big(f_1(\phi(t))\cdot x_1'(t)+\dots+f_k(\phi(t))\cdot x_k'(t)\Big)\ dt\]
\end{statement}
\begin{proof} Перейдем и интегралу римана для каждой координаты:
    \[\int\limits_{\alpha}^{\beta}f_i(\phi(t))\ dx_i(t)=\int\limits_{\alpha}^{\beta}f_i(\phi(t))\cdot x'_i(t)\ dt\]
    при этом
    \[\int\limits_{\gamma}f_1\ dx_1+\dots+f_k\ dx_k=\sum\limits_{i=1}^{k}\int\limits_{\alpha}^{\beta}f_i(\phi(t))\ dx_i(t)=\int\limits_{\alpha}^{\beta}\left(\sum\limits_{i=1}^{k}f_i(\phi(t))\cdot x'_i(t)\right)dt\]
\end{proof}
\begin{statement}
    При смене ориентации кривой $\gamma$, интеграл второго рода меняет знак.
\end{statement}
\begin{proof}
    Рассмотрим интегральную сумму
    \[\sigma(f,T,H)=\sum\limits_{j=1}^{n}\sum\limits_{i=1}^{k}f_i(\phi(\xi_j)) (x_i(t_j)-x_i(t_{j-1}))\]
    При смене ориентации кривой, направление движения по ней меняется на противоположное. Если проходить кривую в обратном направлении то на каждом участке разбиения, начальная и конечные точки меняются местами:
    \[\Delta_{ij}'=x_i(t_{j-1})-x_i(t_j)=-(x_i(t_j)-x_i(t_{j-1}))\]
    Отсюда
    \[\sigma(f,T',H')=-\sum\limits_{j=1}^{n}\sum\limits_{i=1}^{k}f_i(\phi(\xi'_j)) (x_i(t_j)-x_i(t_{j-1}))\]
    Устремим $d(T')\to 0,\ d(T)\to 0$ и получим утверждение теоремы.
\end{proof}
\begin{statement}
    Пусть $\gamma=\gamma_1\sqcup\gamma_2$ и существуют интегралы
    \[\int\limits_{\gamma}f_1\ dx_1+\dots+f_k\ dx_k,\ \int\limits_{\gamma_1}f_1\ dx_1+\dots+f_k\ dx_k,\ \int\limits_{\gamma_2}f_1\ dx_1+\dots+f_k\ dx_k\]
    Тогда 
    \[\int\limits_{\gamma}f_1\ dx_1+\dots+f_k\ dx_k=\int\limits_{\gamma_1}f_1\ dx_1+\dots+f_k\ dx_k+\int\limits_{\gamma_2}f_1\ dx_1+\dots+f_k\ dx_k\]
\end{statement}
\begin{proof}
    По свойству аддитивности для интеграла Римана-Стилтьеса по каждой координат:
    \[\int\limits_{\gamma}f_i\ dx_i=\int\limits_{\gamma_1}f_i\ dx_i+\int\limits_{\gamma_2}f_i\ dx_i\]
    Значит
    \begin{multline*}
        \int\limits_{\gamma}\sum\limits_{i=1}^{k}f_i\ dx_i=\sum\limits_{i=1}^{k}\int\limits_{\gamma}f_i\ dx_i=\sum\limits_{i=1}^{k}\left(\ \int\limits_{\gamma_1}f_i\ dx_i+\int\limits_{\gamma_2}f_i\ dx_i\right)=\\
        =\sum\limits_{i=1}^{k}\int\limits_{\gamma_1}f_i\ dx_i+\sum\limits_{i=1}^{k}\int\limits_{\gamma_2}f_i\ dx_i=\int\limits_{\gamma_1}\sum\limits_{i=1}^{k}f_i\ dx_i+\int\limits_{\gamma_2}\sum\limits_{i=1}^{k}f_i\ dx_i
    \end{multline*}
\end{proof}
\begin{statement}
    Пусть $\gamma\in \mathcal{C}^1$ --- кривая, $f_1,\dots,f_k$ --- непрерывны на $\gamma$. Обозначим $\vec{F}=(f_1,\dots,f_k)$. Тогда
    \[\int\limits_{\gamma}f_1\ dx_1+\dots+f_k\ dx_k=\int\limits_{\gamma}\vec{F}\cdot (x_1'(t),\dots,x_k'(t))\ ds\]
\end{statement}
\begin{comm}
    Пусть $f:\gamma\to \mathbb{C}$ --- комплекснозначная функция, определенная на кривой $\gamma$. Функцию $f$ можно представить в виде
    \[f(x,y)=u(x,y)+i\cdot v(x,y)\]
    где $u(x,y)$ и $v(x,y)$ --- вещественнозначные функции на кривой $\gamma$. Определим криволинейный интерал первого рода рода от комплекснозначной функции: 
    \[\int\limits_{\gamma}f(z)\ |dz|:=\int\limits_{\gamma}u(x,y)\ ds+i\cdot \int\limits_{\gamma} v(x,y)\ ds\]
    Заметим, что $dz=dx+i dy$. Значит
    \begin{multline*}
        f(z)\ dz=(u+iv)(dx+idy)=u\ dx+iu\ dy+iv\ dx+i^2v\ dy=\\
        =(u\ dx-v\ dy)+i(v\ dx+u\ dy)
    \end{multline*}
    Тогда интеграл второго рода:
    \[\int\limits_{\gamma}f(z)\ dz:=\int\limits_{\gamma}u\ dx-v\ dy+\int\limits_{\gamma}v\ dx+u\ dy\]
    Таким образом, криволинейный интеграл второго рода находит широкое применение в комплексном анализе.
\end{comm}
\newpage